\section*{Как читать эту статью}
\addcontentsline{toc}{section}{Как читать эту статью}
\label{sec:reading_guide}

Статья длинная и устроена послойно. Разные читатели приходят с разным интересом; ниже --- короткая карта того, какие разделы наиболее нагружены содержанием для каждой аудитории, и как их можно читать не подряд, а «по своему маршруту».

\begin{center}
\begin{tabular}{p{0.33\linewidth}p{0.60\linewidth}}
\toprule
\textbf{Если вы приходите из\dots} & \textbf{Минимальный маршрут чтения} \\
\midrule
\textbf{клинической практики} & §\ref{subsec:architecture} (архитектура сцепления, первичная/вторичная психопатия), §\ref{subsec:integrated_survivor} (интегрированный выживший), §\ref{sec:limitations} (что модель не поддерживает), §\ref{sec:falsification} C.1--C.2 \\
\midrule
\textbf{психометрии} & §\ref{sec:methods}, §\ref{sec:results} 3.1--3.10, §\ref{sec:falsification} A, §\ref{sec:limitations}.3 (статистические ограничения) \\
\midrule
\textbf{созерцательной или психодинамической традиции} & §\ref{subsec:A_as_field_condition} (A как условие поля), §\ref{subsec:positioning} (позиционирование против Уилбера), §\ref{sec:falsification} D.1 \\
\midrule
\textbf{философии морали и юриспруденции} & §\ref{subsec:architecture} (архитектура vs дефицит), §\ref{sec:falsification} D.2 (P как ортогональный моральный слой) \\
\midrule
\textbf{методологии и open science} & §\ref{sec:methods}~--~\ref{subsec:analytic_primer} (праймер), §\ref{sec:falsification} (фальсификационные обязательства), §\ref{sec:limitations} (self-selection, самоотчёт, cross-sectional) \\
\midrule
\textbf{теории и программы} & \S\ref{sec:discussion}, §\ref{sec:future}, \texttt{PROGRAM\_STRUCTURE.md} \\
\bottomrule
\end{tabular}
\end{center}

\paragraph{Короткий справочник шести факторов.} Чтобы не возвращаться каждый раз к определениям, ниже --- одностраничная сводка того, что измеряет каждый из шести публикуемых факторов, его надёжность на наших данных и концептуальный корень.

\begin{center}
\small
\begin{tabular}{p{0.08\linewidth}p{0.22\linewidth}cp{0.23\linewidth}p{0.30\linewidth}}
\toprule
\textbf{Фактор} & \textbf{Что меряет} & \boldmath$\omega$ & \textbf{Wave-1 прародитель} & \textbf{Ключевые линии} \\
\midrule
ID & Внутренний диссонанс: чувствование запущено, но не идёт дальше; конфликт между affective impulse и блокировкой & 0.91 & IDS (Internal Dissonance Scale) + S, D, AIS, F как surface-манифестации & \citet{ferenczi1933, klein1946, fonagy2002}; защиты как сеть \citep{digiuseppe2024} \\
\midrule
CA & Детское отчуждение: ранняя реляционная история, в которой внутренние состояния ребёнка не встречали адекватного отклика & 0.89 & ACE + CA, объединённые в один латент & \citet{bowlby1969, main1986, fonagy2002}; раннее пренебрежение как предиктор защитной архитектуры \\
\midrule
P & Моральный нарратив «это разрешено мне»: слой разрешения на использование других, структурно отделимый от аффективной архитектуры & 0.93 & P wave-1 (без изменения) & \citet{cleckley1941, hare1991, blair2013, baroncohen2011}; LSRP \citep{levenson1995} \\
\midrule
M & Маскирование: регуляторный слой между внутренним состоянием и внешним предъявлением. \textbf{Слабо дифференцирован} ($r_{\mathrm{lat}}(M, P) = +0.82$); самостоятельных выводов не несёт & 0.55 & MS wave-1 & Impression management; редизайн в волне~2 \\
\midrule
\EAF{} & Аффективная эмпатия: автоматический резонанс с чужим состоянием & 0.90 & E (общая), post-hoc расщеплена & \citet{davis1983, shamaytsoory2011, decety2011} \\
\midrule
\Ecog{} & Когнитивная эмпатия: распознавание и понимание чужих ментальных состояний без обязательного аффективного отклика & 0.89 & E (общая), post-hoc расщеплена & \citet{baroncohen2011, shamaytsoory2011} \\
\bottomrule
\end{tabular}
\end{center}

Плюс седьмой конструкт --- \textbf{внимание (A)}, работающий как формативный композит (не рефлексивный латент), инструментированный через 31 айтем из 9 шкал с теоретически взвешенными коэффициентами. A --- не содержание психики, а условие её организации (§\ref{subsec:A_as_field_condition}). В волне~2 планируется отдельная A-шкала.

\paragraph{Архитектурная триада D--P--A.} В рабочем языке статьи и программы мы используем короткое обозначение \textbf{D--P--A} для трёх содержательно неэквивалентных осей архитектуры. Обозначение исторически досталось от ранних версий инструмента, где \textbf{D} обозначало общий \emph{Defense}-бандл --- агрегат всех защитных подшкал (подавление, диссоциация, аффективная изоляция, уплощение, внутренний диссонанс), а не какую-либо одну из них в отдельности. В текущей модели функциональную роль этого бандла выполняет латент \textbf{ID}: тот же самый набор защит, поглощённый в единый психометрический конструкт на основании эмпирически высокой взаимной корреляции подшкал ($r = 0.77$--$0.94$) и сходимости с результатами сетевого анализа защит \citep{digiuseppe2024}. Сохраняем трёхбуквенный ярлык потому, что он содержательно удобен --- описывает три функционально и онтологически разные слоя одной системы:

\begin{center}
\begin{tabular}{clp{0.62\linewidth}}
\toprule
\textbf{Ось} & \textbf{Латент} & \textbf{Функция} \\
\midrule
\textbf{D} & ID & \textbf{Пассивное} не-присоединение. Аффективный сигнал запущен, но дальше не идёт: блокировка, подавление, диссоциация, аффективная изоляция, аффективная уплощённость как поверхностные манифестации одного поля. Этиологически --- через CA (детское отчуждение). Эго-дистонный: субъект \emph{хочет} чувствовать, но не может. \\
\midrule
\textbf{P} & P & \textbf{Активное} не-присоединение. Аффективный резонанс структурно доступен, но запускается условно; сопровождается моральным нарративом «мне это разрешено». Этиологически --- ортогонален CA, то есть не является следствием детской истории на популяционном уровне. Эго-синтонный: субъект \emph{выбирает} не чувствовать. \\
\midrule
\textbf{A} & A (формат.) & \textbf{Условие поля}, не содержание. Наблюдающее отношение, свидетельствование, мета-осознавание. Онтологически асимметрично D и P: не ещё одна функция, а режим, в котором существуют остальные функции. В наших данных --- модератор обеих стадий канала CA$\to$ID$\to$\EAF{} (в узкой спецификации) и главный позитивный предиктор \EAF{} в объединённой модели. \\
\bottomrule
\end{tabular}
\end{center}

Соответственно, когда в тексте статьи говорится о «D-канале» или «D-паттерне», имеется в виду пассивно-блокирующая архитектура с латентом ID как психометрическим якорем; «P-канал» --- активно-нарративная архитектура; «A-слой» --- свидетельствующий мета-уровень. Отсюда же вытекают описания типов: первичная психопатия --- \textbf{P без D}; вторичная психопатия --- \textbf{P с D}; интегрированный выживший --- \textbf{D с A}; базовый эмпатический паттерн --- \textbf{ни D ни P}.

\paragraph{Методологический глоссарий.} Чтобы не переводить технические термины по ходу чтения, ниже --- одностраничная сводка русско-английских соответствий для методологических понятий, используемых в тексте. Там, где английский термин сохранён (аббревиатуры, имена), это отмечено.

\begin{center}
\small
\begin{tabular}{p{0.38\linewidth}p{0.56\linewidth}}
\toprule
\textbf{Русский термин в тексте} & \textbf{Английский эквивалент / уточнение} \\
\midrule
подтверждающий факторный анализ, CFA & confirmatory factor analysis (CFA) \\
эксплоративный факторный анализ & exploratory factor analysis (EFA) \\
косоугольная / ортогональная CFA & oblique / orthogonal CFA \\
латентный фактор, латентная переменная & latent factor / latent variable \\
факторная нагрузка & factor loading \\
факторные оценки & factor scores \\
индексы соответствия & fit indices \\
индекс сравнительного соответствия & Comparative Fit Index (CFI) \\
средняя ошибка аппроксимации & Root Mean Square Error of Approximation (RMSEA) \\
бифакторная модель & bifactor model \\
доля объяснённой общей дисперсии & Explained Common Variance (ECV) \\
иерархическая омега & hierarchical $\omega$ \\
главный эффект / эффект взаимодействия & main effect / interaction effect \\
модерация & moderation \\
медиация & mediation \\
параллельная медиация & parallel mediation \\
непрямой эффект / путь~$a$, путь~$b$, прямой путь $c'$ & indirect effect / $a$-path, $b$-path, direct path $c'$ \\
модерированная медиация & moderated mediation \\
двойно-модерированная медиация & moderated-moderated mediation \\
индекс двойно-модерированной медиации & Index of Moderated-Moderated Mediation (\IMM) \\
объединённое структурное уравнение & unified SEM \\
бутстрэп / бутстрэповский доверительный интервал & bootstrap / bootstrap confidence interval \\
апостериорное распределение, байесовский постериор & Bayesian posterior \\
приор, априорное распределение & prior \\
95\%-ный интервал наивысшей плотности & 95\% highest density interval (HDI) \\
масштабирующий множитель, коррекция Саторра--Бентлера & Satorra--Bentler scaling factor \\
рефлексивный / формативный конструкт & reflective / formative construct \\
контроль переменной, частная корреляция & partialling / partial correlation \\
робастность, проверка устойчивости & robustness check \\
метод расщеплённой половины & split-half method \\
инвариантность измерения & measurement invariance \\
конфигуральная / метрическая / скалярная инвариантность & configural / metric / scalar invariance \\
характеристическая кривая, площадь под кривой & ROC curve / area under the curve (AUC) \\
\bottomrule
\end{tabular}
\end{center}

\paragraph{Что статья не делает.} Чтобы избежать разочарования от неверно направленных ожиданий: это \emph{не} longitudinal study, \emph{не} клинически валидированная шкала с cut-off scores, \emph{не} кросс-культурная валидация, \emph{не} попытка заместить DSM. Это --- первое эмпирическое валидирование шестифакторной архитектуры на русскоязычной онлайн-выборке, с эксплицитным списком того, что следующие волны программы должны проверить (§\ref{sec:falsification}, §\ref{sec:future}). Cross-sectional, self-report, моноязычный.

\clearpage
